% ============================================================
% styles2.tex
% Encabezado LaTeX para los TALLERES (en clase y en casa) del curso
% Estadistica Multivariada - Programa de Biologia, Universidad del Magdalena
% Prof. Javier Rodriguez Barrios
%
% Se incluye en cada taller .qmd asi:
%   format:
%     pdf:
%       include-in-header: styles2.tex
%
% IMPORTANTE: el proyecto usa documentclass: scrartcl (KOMA-Script), NO la
% clase "article" estandar. Por eso este archivo usa SOLO mecanismos nativos
% de KOMA-Script para encabezados y estilos de titulo (\addtokomafont,
% scrlayer-scrpage), y evita a proposito el paquete "titlesec" (rompe la
% compilacion con scrartcl: error "Argument of \paragraph has an extra }")
% y el paquete "tocloft" (genera advertencias/conflictos con el sistema de
% tabla de contenido propio de KOMA-Script).
%
% Objetivo: una plantilla SENCILLA y didactica (no editorial tipo libro),
% facil de leer para los estudiantes y facil de replicar por ellos mismos
% en sus propios informes/manuscritos en RStudio.
% ============================================================

% ---------- 1) Tipografia segura ----------
% lmodern solo tiene sentido con pdflatex (motor de 8-bits); con xelatex
% (el motor que usa este proyecto) las fuentes las maneja fontspec, asi que
% lmodern se carga condicionalmente para no interferir.
\usepackage{iftex}
\ifPDFTeX
  \usepackage{lmodern}
\fi

% ---------- 2) BLOQUES DE CODIGO: nunca se salen del margen ----------
% fvextra controla los entornos "Verbatim" que usa Pandoc/Quarto para los
% chunks de codigo (Highlighting) y para las salidas de R.
\usepackage{fvextra}
\fvset{
  breaklines    = true,  % permite cortar la linea
  breakanywhere = true,  % corta aunque no haya espacios (rutas, strings largos, etc.)
  fontsize      = \small,
  breaksymbolleft = {}, breaksymbolright = {}   % sin flechas de continuacion: mas limpio para los estudiantes
}
% Redefinicion explicita del entorno que usa Pandoc/Quarto para el codigo con
% resaltado de sintaxis (chunks de R). "commandchars" es INDISPENSABLE aqui:
% permite que \NormalTok{}, \FunctionTok{}, \StringTok{}, etc. (los comandos
% de color que Pandoc inserta token por token) se sigan interpretando como
% LaTeX real dentro del entorno verbatim; sin esta opcion se verian como texto
% literal en vez de como codigo coloreado.
\DefineVerbatimEnvironment{Highlighting}{Verbatim}{
  breaklines, breakanywhere, fontsize=\small, breaksymbolleft={}, breaksymbolright={}, commandchars=\\\{\}
}
% Redefinicion del entorno "verbatim" comun (lo usa, por ejemplo, la salida de
% consola de R que no lleva resaltado de sintaxis). Por defecto este entorno
% NO corta linea, asi que sin esto una salida larga sigue desbordando el margen.
\DefineVerbatimEnvironment{verbatim}{Verbatim}{
  breaklines, breakanywhere, fontsize=\small, breaksymbolleft={}, breaksymbolright={}
}

% Fondo gris suave para los chunks (entorno "Shaded" que ya define Pandoc;
% aqui solo se le cambia el color de fondo)
\usepackage{framed}
\renewenvironment{Shaded}{%
  \definecolor{shadecolor}{HTML}{F5F7FA}%
  \begin{snugshade}%
}{\end{snugshade}}

% ---------- 3) Paleta de color del curso ("Analitico") ----------
\usepackage{xcolor}
\definecolor{azulcurso}{HTML}{0C447C}   % azul profundo
\definecolor{azulclaro}{HTML}{85B7EB}   % azul cielo (linea de encabezado)
\definecolor{ambarcurso}{HTML}{EF9F27}  % ambar de acento

% ---------- 4) Enlaces y referencias cruzadas SIEMPRE en azul ----------
% (Aqui no aplican normas editoriales de libro: todo enlace, cita o
%  referencia cruzada se resalta en azul para que el estudiante la
%  identifique de inmediato como algo "clicable"/consultable).
% NOTA IMPORTANTE: como el YAML del taller trae "colorlinks: true", Pandoc
% inserta su PROPIO \hypersetup (con enlaces en color "Maroon"/rojizo) justo
% DESPUES de este archivo, lo que pisaria estos colores si se definieran aqui
% directamente. Por eso el \hypersetup real se aplica con \AtBeginDocument,
% que se ejecuta al iniciar el documento, es decir, DESPUES de que Pandoc ya
% fijo sus propios colores por defecto -- asi estos colores de azules del
% curso son siempre los que quedan.
\usepackage{hyperref}
\AtBeginDocument{%
  \hypersetup{
    colorlinks = true,
    linkcolor  = azulcurso,   % referencias cruzadas internas (secciones, figuras, tablas)
    citecolor  = azulcurso,   % citas bibliograficas (biblio.bib / apa.csl)
    urlcolor   = azulcurso,   % enlaces web (p. ej. al repositorio de talleres)
    filecolor  = azulcurso,
    linktoc    = all,         % toda la tabla de contenido queda enlazada
    breaklinks = true
  }%
}

% ---------- 5) Titulos de seccion: color del curso, con el mecanismo NATIVO
%               de KOMA-Script (\addtokomafont), sin tocar tamanos/espaciados
%               que ya trae bien calculados scrartcl con DIV=11. ----------
\addtokomafont{section}{\color{azulcurso}}
\addtokomafont{subsection}{\color{azulcurso}}
\addtokomafont{subsubsection}{\color{ambarcurso}}
% El taller usa "block-headings: true" (definido a nivel de proyecto), que
% hace que Pandoc traduzca "####" a \subparagraph (no a \paragraph) para que
% se vea como un encabezado de bloque real en vez de texto "run-in". Por eso
% se recolorea \subparagraph aqui, con \RedeclareSectionCommand (la
% herramienta nativa de KOMA-Script para esto, ya que \addtokomafont no
% tiene efecto visible en este nivel).
\RedeclareSectionCommand[font=\normalfont\bfseries\color{ambarcurso}]{subparagraph}

% ---------- 6) Encabezado y pie de pagina simples (KOMA-Script nativo:
%               scrlayer-scrpage, en vez de fancyhdr) ----------
% \tallertipo se puede sobrescribir en un taller puntual asi (en el YAML del .qmd):
%   header-includes:
%     - \renewcommand{\tallertipo}{Taller de entrenamiento}
\providecommand{\tallertipo}{Taller en clase}

\usepackage{scrlayer-scrpage}
\clearpairofpagestyles
\ihead{\textcolor{azulcurso}{\small\bfseries \tallertipo\ -- Estad\'istica Multivariada}}
\ohead{\small\pagemark}
\cfoot{\small\color{gray}Universidad del Magdalena -- Programa de Biolog\'ia}
\KOMAoptions{headsepline = true, footsepline = false}
\setkomafont{headsepline}{\color{azulclaro}}
\pagestyle{scrheadings}
% \maketitle usa por defecto el estilo de pagina "plain" (sin encabezado) para
% la portada. Se redirige a "scrheadings" para que la portada tambien lleve
% el mismo encabezado que el resto del taller.
\renewcommand*{\titlepagestyle}{scrheadings}

% ---------- 7) Tablas (kableExtra) legibles y sin desbordes ----------
\usepackage{longtable}
\usepackage{booktabs}
\usepackage{array}
